\startcomponent ch04
\product fduthesis

\chapter{引用与参考文献}

\section{参考文献管理}

本模板使用 \ConTeXt\ 内置的参考文献系统。
参考文献条目存放在 \type{references.bib} 文件中，采用标准 BibTeX 格式。

在主文件中通过以下命令加载文献数据库：

\starttyping
\usebtxdataset[references.bib]
\stoptyping

\section{引用方式}

使用 \type{\cite[key]} 命令引用文献。以下是几种引用示例。

Vaswani 等人提出了 Transformer 架构
\cite[vaswani2017attention]，
该架构完全基于自注意力机制，
摒弃了传统的循环和卷积结构。
在此基础上，Devlin 等人提出了 BERT 模型
\cite[devlin2019bert]，
通过大规模预训练实现了多项自然语言处理任务的突破。
关于深度学习的基础理论，
可参阅 Goodfellow 等人的著作 \cite[goodfellow2016deep]。

近年来，大语言模型的研究取得了重大进展。
Brown 等人提出的 GPT-3 \cite[brown2020gpt3]
展示了少样本学习的强大能力。
Liu 等人的 RoBERTa \cite[liu2019roberta]
通过优化预训练策略进一步提升了性能。

\section{文献格式}

根据复旦大学论文规范，
参考文献编排格式须符合国家标准
GB/T 7714—2015
《信息与文献：参考文献著录规则》。
\ConTeXt\ 内置的 APA 风格与该标准有一定差异，
用户可根据需要进行微调。

常见的参考文献类型及其 BibTeX 格式如下：

\starttyping
% 期刊论文
@article{key,
  author  = {作者},
  title   = {论文标题},
  journal = {期刊名称},
  year    = {2024},
  volume  = {1},
  pages   = {1--10},
}

% 会议论文
@inproceedings{key,
  author    = {作者},
  title     = {论文标题},
  booktitle = {会议名称},
  year      = {2024},
  pages     = {100--110},
}

% 图书
@book{key,
  author    = {作者},
  title     = {书名},
  publisher = {出版社},
  year      = {2024},
}
\stoptyping

\section{交叉引用}

除了引用参考文献，还可以交叉引用公式、表格、图片和章节：

\startitemize
  \item 引用公式：\type{\in{公式}{~}[eq:crossentropy]}
        \rightarrow\ \in{公式}{~}[eq:crossentropy]
  \item 引用表格：\type{\in{表}{~}[tab:comparison]}
        \rightarrow\ \in{表}{~}[tab:comparison]
  \item 引用章节：\type{\in{第}{章}[ch01]}
\stopitemize

\stopcomponent
